\documentclass[a4paper,11pt]{article}
\usepackage[utf8]{inputenc}
\usepackage[T1]{fontenc}
\usepackage[french]{babel}
\usepackage{lmodern}
\usepackage{amsmath,amssymb,amsthm}
\usepackage{mathrsfs}
\usepackage{enumitem}
\usepackage{multicol}

\setlist[enumerate]{label=\textbf{\textcolor{Couleur8}{\arabic*.}}, left=1em}

% Si vous voulez aussi modifier les listes enumerate imbriquées :
\setlist[enumerate,2]{label=\textcolor{Couleur8}{\alph*)}}
\setlist[enumerate,3]{label=\textcolor{Couleur8}{\roman*)}}
\setlist[enumerate,4]{label=\textcolor{Couleur8}{\Alph*)}}
\usepackage{graphicx}
\usepackage{xcolor}
\usepackage{tcolorbox}
\usepackage{geometry}
\usepackage{fancyhdr}
\usepackage{titlesec}
\usepackage{cancel}
\usepackage{ifthen}
\usepackage{tikz}
\usetikzlibrary{arrows.meta}
\usepackage{tkz-tab}
\usepackage{eurosym}

% OPTION PROF/ÉLÈVE
% Décommentez UNE des deux lignes suivantes :
\newboolean{prof}\setboolean{prof}{true}   % Version professeur (avec corrections des devoirs)
%\newboolean{prof}\setboolean{prof}{false}  % Version élève (sans corrections des devoirs)

\definecolor{Couleur1}{HTML}{4A5D7A} %Th
\definecolor{Couleur2}{HTML}{A5B5D6} %Prop
\definecolor{Couleur3}{HTML}{A8C68A} %Méthode
\definecolor{Couleur6}{HTML}{8A9CC1} %Def
\definecolor{Couleur7}{HTML}{E6B459} %Rq Voc
\definecolor{Couleur8}{HTML}{D36740} %Titres
\definecolor{correction}{HTML}{D36740} % Couleur pour les corrections
\definecolor{coopmathscorr}{HTML}{F15929} % Couleur corrections coopmaths

% Configuration des marges
\geometry{
	top=2cm,
	bottom=2cm,
	inner=1.5cm,
	outer=1.5cm,
}

% Police
\renewcommand{\familydefault}{\sfdefault}

% Compteurs
\newcounter{seancenum}
\newcounter{seancenumD}
\setcounter{seancenum}{1}
\setcounter{seancenumD}{1}

% Configuration des en-têtes et pieds de page
\pagestyle{fancy}
\fancyhf{}
\ifthenelse{\boolean{prof}}{
    \fancyhead[L]{\textcolor{Couleur8}{\textbf{Corrections Automatismes - Classe de seconde - Version Professeur}}}
}{
    \fancyhead[L]{\textcolor{Couleur8}{\textbf{Corrections Automatismes - Séance 9 - Classe de seconde}}}
}
\fancyhead[R]{\textcolor{Couleur8}{\textbf{2025-2026}}}
\fancyfoot[L]{\textcolor{Couleur8}{Mme Bertail et Mme Pinto}}
\fancyfoot[R]{\textcolor{Couleur8}{\thepage}}
\renewcommand{\headrulewidth}{1pt}
\renewcommand{\headrule}{\hbox to\headwidth{\color{Couleur8}\leaders\hrule height \headrulewidth\hfill}}

% Environnements personnalisés
\newtcolorbox{seance}[1]{
    colback=Couleur6!10,
    colframe=Couleur1!65,
    fonttitle=\bfseries\large,
    title=Correction Séance \theseancenum #1,
    left=5mm,
    right=5mm,
    top=3mm,
    bottom=3mm,
    arc=0mm,
    after=\stepcounter{seancenum}
}

\newtcolorbox{seanceD}[1]{
    colback=Couleur6!5,
    colframe=Couleur6!45,
    fonttitle=\bfseries\large,
    title=Correction Devoirs - Séance \theseancenumD #1,
    left=5mm,
    right=5mm,
    top=3mm,
    bottom=3mm,
    arc=0mm,
    after=\stepcounter{seancenumD}
}

\begin{document}

\setcounter{seancenum}{9}
\newpage
\begin{seance}{} %9
\begin{enumerate}
 \item Diminuer de $20\,\%$ revient à multiplier par $0{,}8$ et augmenter de $20\,\%$ revient à multiplier par $1{,}2$.\\
  Globalement cela revient donc à multiplier par $0{,}8\times 1{,}2=0{,}96$.\\
  Multiplier par $0{,}96$ revient à multiplier par  $1-0{,}04$. \\
          Le taux d'évolution global est donc : ${\color[HTML]{f15929}\boldsymbol{-4\,}} \%$.


\item Le coefficient multiplicateur associé à une baisse de $75\,\%$ est $1-0{,}75=0{,}25$.\\
    Comme $0{,}25\times 4=1$, il faut donc multiplier par $4$ (coefficient multiplicateur réciproque) pour revenir au prix initial.\\
      Un coefficient multiplicateur de $4$ correspond à un taux d'évolution de $300\,\%$.\\
      On en déduit qu'il faut une augmentation de ${\color[HTML]{f15929}\boldsymbol{300\,\%}}$.\\La bonne réponse est la réponse {\bfseries \color[HTML]{f15929}B}.


\item Il s'agit d'un problème additif. Il va être nécessaire de réduire les fractions au même dénominateur pour les additionner, les soustraire ou les comparer.\\Réduisons les fractions de l'énoncé au même dénominateur :  $\dfrac{3}{10}$ $= \dfrac{12}{40}$ ; $\dfrac{1}{4}$ $= \dfrac{10}{40}$ et $\dfrac{7}{40}$ .\\Calculons d'abord la fraction du jardin occupée par une serre servant aux semis : \\$1-\dfrac{3}{10}-\dfrac{1}{4}-\dfrac{7}{40} = \dfrac{40}{40}-\dfrac{12}{40}-\dfrac{10}{40}-\dfrac{7}{40} = \dfrac{40-12-10-7}{40} = \dfrac{11}{40}$\\Le jardin est donc occupé de la façon suivante : $\dfrac{3}{10}$ par la culture des légumes, $\dfrac{1}{4}$ par la culture des plantes aromatiques, $\dfrac{7}{40}$ par une serre servant aux semis et $\dfrac{11}{40}$ par la culture des fraisiers.\\ Avec les mêmes dénominateurs pour pouvoir comparer, le jardin est donc occupé de la façon suivante : $\dfrac{12}{40}$ par la culture des légumes, $\dfrac{10}{40}$ par la culture des plantes aromatiques, $\dfrac{7}{40}$ par une serre servant aux semis et $\dfrac{11}{40}$ par la culture des fraisiers.\\Nous pouvons alors ranger ces fractions dans l'ordre croissant : $\dfrac{7}{40}$, $\dfrac{10}{40}$, $\dfrac{11}{40}$, $\dfrac{12}{40}$.\\Enfin, nous pouvons ranger les fractions de l'énoncé et la fraction calculée dans l'ordre croissant : ${\color[HTML]{f15929}\boldsymbol{\dfrac{7}{40}}} < {\color[HTML]{f15929}\boldsymbol{\dfrac{1}{4}}} < {\color[HTML]{f15929}\boldsymbol{\dfrac{11}{40}}} < {\color[HTML]{f15929}\boldsymbol{\dfrac{3}{10}}}$.\\ 
                      C'est donc par {\bfseries \color[HTML]{f15929}la culture des légumes} que le jardin est le plus occupé.


\item
            Pour $x\in \mathbb{R}\smallsetminus\left\{\dfrac{-1}{2}\,;\,\dfrac{1}{3}\right\}$, \\
            $\begin{aligned}
            \dfrac{3}{9x-3}+\dfrac{1}{2x+1}
            &= \dfrac{3(2x+1)}{(9x-3)(2x+1)}+\dfrac{1(9x-3)}{(9x-3)(2x+1)}\\
            &=\dfrac{3(2x+1)+(9x-3)}{(9x-3)(2x+1)}\\
            &=\dfrac{15x}{(9x-3)(2x+1)}
            \end{aligned}$


\item $4(-2x-7)=-9x-3$\\On développe le membre de gauche.\\$-8x-28=-9x-3$\\On ajoute $9x$ aux deux membres.\\$-8x-28{\color[HTML]{216D9A}\boldsymbol{+9x}}=-9x-3{\color[HTML]{216D9A}\boldsymbol{+9x}}$\\$x-28=-3$\\On ajoute $28$ aux deux membres.\\$x-28{\color[HTML]{216D9A}\boldsymbol{+28}}=-3{\color[HTML]{216D9A}\boldsymbol{+28}}$\\$x=25$\\On divise les deux membres par $1$.\\$x{\color[HTML]{216D9A}\boldsymbol{\div1}}=25{\color[HTML]{216D9A}\boldsymbol{\div1}}$\\$x=25$\\La solution de l'équation $4(-2x-7)=-9x-3$ est ${\color[HTML]{f15929}\boldsymbol{25}}$.

\end{enumerate}
\end{seance}

\end{document}